% Imporditud: tools/import_eksperimendid.py
% Allikas: Eesti füüsikaolümpiaadi eksperimendi ülesannete
%   kogu (Rael Kalda, 2023), ülesanne Ü116, lahendus L116.
\setAuthor{EFO žürii}
\setRound{piirkonnavoor}
\setYear{2017}
\setNumber{P E2}
\setDifficulty{3}
\setTopic{termodünaamika}
\setVahendid{kalorimeeter, termomeeter, joonlaud, vesi, lumi}
\setPunktid{12}
\setAllikas{Eksperimendi kogu Ü116, lahendus lk 91}
\setLahendusseis{toores}

\prob{Kalorimeeter}
Leidke kalorimeetri soojusmahtuvus Ckal.
\textit{Vihje:} Kalorimeetri soojusmahtuvus on kalorimeetri topsi erisoojuse ja massi korrutis Ckal = ckalmkal. Vee erisoojus cv = 4200 J/(kg $\cdot ^\circ$C), jää sulamissoojus $\lambda$ = 330 kJ/kg, vee tihedus $\rho$ = 1,0 g/cm$^3$.

\hint

\solu
Täidame umbes pool kalorimeetrist toasooja veega (temperatuur $T_0$). Vee $\pi$ru(ud2m)2ahla. V saame mõõta kalorimeetri diameetri d ning vee kõrguse h kaudu. V =

Kuna vee tihedus $\rho$ = 1 g/cm$^3$, siis saame kalorimeetris oleva vee massiks

( d )2 mvesi = $\pi$ 2 h$\rho$.

Lisame kalorimeetrisse lund nii, et vee temperatuur langeks alla 10 $^\circ$C ning määrame vee lõpptemperatuuri $T_1$. Väike temperatuuri muutus annab väga ebatäpse tulemuse. Lisatud lume massi saame kätte sulanud vee ruumala kaudu. Mõõdame uuesti veetaseme $h_2$ kalorimeetris ning arvutame veetaseme tõusu $\Delta$h = $h_2$ $-$ h. Seega lisatud lume mass oli

( d )2 mlumi = $\pi$ 2 ($h_2$ $-$ h)$\rho$.

Lume sulamiseks vajaminev soojushulk

$Q_1$ = $\lambda$mlumi Tekkinud vee temperatuuri tõstmiseks temperatuurini $T_1$ vajaminev soojushulk

$Q_2$ = cmlumi($T_1$ $-$ 0 $^\circ$C). saadakse kalorimeetris olnud vee jahtumise

$Q_3$ = cmvesi($T_0$ $-$ $T_1$) ning kalorimeetri enda jahtumise arvel

Qkal = ckalmkal($T_0$ $-$ $T_1$).

Teades, et $Q_1$ + $Q_2$ = $Q_3$ + $Q_4$ ning arvutades välja $Q_1$, $Q_2$ ja Qkal, saame leida kalorimeetri soojusmahtuvuse Ckal

Ckal = $Q_1$ + $Q_2$ $-$ $Q_3$ . ($T_0$ $-$ $T_1$)

Hindamine: Kalorimeetris oleva vee ruumala/massi mõõtmine --- 2 p. Kalorimeetrisse lisatud lume massi mõõtmine sulanud vee järgi --- 2 p. Vee alg- ja lõpptemperatuuri mõõtmine --- 2 p. Soojushulkade $Q_1$, $Q_2$ ja $Q_3$ arvutamine --- 3 p. Kalorimeetri soojusmahtuvuse Ckal arvutamine --- 2 p. Kogu katse tegemine vähemalt kaks korda --- 1 p.

Märkus: Kui katses kasutatakse lund, siis mõjutab tulemust ka see, et kasutatav lumi on märg, mistõttu lume sulamiseks vajaminevat soojushulka $Q_1$ ei saa täpselt leida (ei ole võimalik hinnata seda, kui palju lumi sisaldab vett). Marja lume mainimine/mitte mainimine punkte ei mõjuta.
\probend
